\documentclass[11pt,letterpaper]{article}

%% === margins ===
\addtolength{\hoffset}{-0.8in} \addtolength{\voffset}{-0.8in}
\addtolength{\textwidth}{1.6in} \addtolength{\textheight}{1.6in}

%% JASA format with 12pt, spacingset = 1.83
%\addtolength{\hoffset}{-0.3in} \addtolength{\voffset}{-1.2in}
%\addtolength{\textwidth}{.6in} \addtolength{\textheight}{2.1in}
\pdfminorversion=4

%% === basic packages ===
\usepackage{lmodern}% http://ctan.org/pkg/lm
\usepackage{latexsym,multirow}
\usepackage{amssymb,amsmath,bm}
\usepackage{bigints}
\usepackage{comment}
\usepackage{bbm} % for a indicator function
\usepackage{graphicx,subfigure}
\usepackage[color,all,import,arrow]{xy}
\usepackage{enumerate}
\usepackage{verbatim}
\usepackage{mathtools}
\usepackage{longtable}
\usepackage{fancyvrb}

% bold symbol in headers
\usepackage{unicode-math}


%% === graphical packages ===
\usepackage{placeins}

%% === bibliography packages ===
%\usepackage{natbib}
\usepackage[natbib=true,minnames=15,maxnames=15,backend=biber,style=authoryear-luh-ipw]{biblatex}
\addbibresource{mybib.bib}
% \bibliographystyle{pa}

%% === hyperref options ===
% \usepackage{color}
\usepackage[pdftex, bookmarksopen=true, bookmarksnumbered=true,
pdfstartview=FitH, breaklinks=true, urlbordercolor={0 1 0}, citebordercolor={0 0 1}]{hyperref}
\usepackage{colortbl}
\usepackage{subfigure}

%% ==clean up hanging first page==
\usepackage{atbegshi}% http://ctan.org/pkg/atbegshi
\AtBeginDocument{\AtBeginShipoutNext{\AtBeginShipoutDiscard}}

% === dcolumn package ===
\usepackage{dcolumn}
\newcolumntype{.}{D{.}{.}{-1}}
\newcolumntype{d}[1]{D{.}{.}{#1}}

% === figure path ===
\graphicspath{{./Figures/}}

% === caption ===
\usepackage{caption}

% == tikzpicture ==
\usepackage[T1]{fontenc}
\usepackage{tikz}
\usepackage{qtree}
\usepackage{comment}
\usetikzlibrary{arrows}
\usetikzlibrary{positioning}
\usetikzlibrary{matrix}
\usetikzlibrary{bayesnet}

% == Algorithm
%\usepackage[ruled,vlined]{algorithm2e}
%FOR ALGORITHM 
\usepackage{algorithmicx}
\usepackage{algorithm} 
\newcommand*\Let[2]{\State #1 $\gets$ #2}
\algrenewcommand\alglinenumber[1]{
    {\sf\footnotesize\addfontfeatures{Colour=888888,Numbers=Monospaced}#1}}
\algrenewcommand\algorithmicrequire{\textbf{Precondition:}}
\algrenewcommand\algorithmicensure{\textbf{Postcondition:}}
\usepackage[noend]{algpseudocode}
%END FOR ALGORITHM 

% === theorem package ===
\usepackage{theorem}
\theoremstyle{plain}
\theoremheaderfont{\scshape}
\newtheorem{assumption}{Assumption}
\newtheorem{example}{Example}
\newtheorem{definition}{Definition}
\newtheorem{corollary}{Corollary}
\newtheorem{property}{Property}
\newtheorem{proposition}{Proposition}
\newtheorem{theorem}{Theorem}
\newtheorem{defi}{Definition}
\newtheorem{thm}{Theorem}
\newtheorem{lemma}{Lemma}

% === some special symbols
\newcommand{\iid}{\stackrel{\rm i.i.d.}{\sim}}
%\newcommand{\indep}{\stackrel{\rm indep.}{\sim}}
\newcommand{\indep}{\mbox{$\perp\!\!\!\perp$}}
\newcommand{\qed}{\hfill \ensuremath{\Box}}
\newcommand{\ind}{\mbox{$\perp\!\!\!\perp$}}
\newcommand{\nind}{\mbox{$\not\perp\!\!\!\perp$}}
\def\independenT#1#2{\mathrel{\rlap{$#1#2$}\mkern2mu{#1#2}}}
\DeclareMathOperator{\sgn}{sgn}
\DeclareMathOperator{\diag}{diag}
\providecommand{\norm}[1]{\lVert#1\rVert}

% ==== rotating package ===
\usepackage{rotating}

% ==== dotted lines in tables ===
\usepackage{arydshln}

% == spacing between sections and subsections
\usepackage[compact]{titlesec}
%\usepackage{times}

% == control space within footnote
\usepackage{setspace}
\usepackage[bottom]{footmisc}
\renewcommand{\footnotelayout}{\setstretch{1.1}}% Footnotes are \setstretch{1.5}
\setlength{\footnotemargin}{1mm}


\allowdisplaybreaks

\newcommand\spacingset[1]{\renewcommand{\baselinestretch}%
{#1}\small\normalsize}

%%%%%%%%%%%%%%%%%%%%%%%%%%%%%%%%%%%%%%%%%%%%%%%%%%%%%%%%%%%%%%%%%%%%%%

%% === submission
\newcommand{\blind}{1}

\def\letas{\mathrel{\mathop{=}\limits^{\triangle}}}

\newcommand*{\QEDB}{\hfill\ensuremath{\square}}
\newcommand{\SimDateOne}{2022DASH09DASH30}
\newcommand{\SimDateTwo}{2022DASH08DASH10}
\newcommand{\red}{\color{red}}
\newcommand{\blue}{\color{blue}}
\newcommand{\Beta}{\textsf{Beta}}
\newcommand{\Binomial}{\text{Binomial}}
\newcommand{\Bern}{\textsf{Bernoulli}}
\newcommand{\Expo}{\textsf{Expo}}
\newcommand{\Pois}{\textsf{Pois}}
\newcommand{\Unif}{\textsf{Uniform}}
\newcommand{\Normal}{\mathcal{N}}
\newcommand{\Gammad}{\textsf{Gamma}}
\newcommand{\logit}{\text{logit}}
\newcommand{\expit}{\text{expit}}
\newcommand{\mexpit}{\text{mexpit}}
\newcommand{\Dir}{\textsf{Dirichlet}}
\newcommand{\Multi}{\textsf{Multi}}
\newcommand{\Cat}{\textsf{Categorical}}
\newcommand{\Neural1}[1]{\textsf{Neural$_1$}\{#1\}}
\newcommand{\Neural2}[1]{\textsf{Neural$_2$}\{#1\}}

\newcommand{\pr}{\text{pr}}
\newcommand{\Var}{\textrm{Var}}
\newcommand{\Cov}{\text{Cov}}
\newcommand{\softmax}{\textrm{softmax}}
\newcommand{\sumN}{\sum_{i=1}^N}
\newcommand{\wt}{\widetilde}

\newcommand{\Prr}{\textrm{Pr}}
\newcommand{\E}{\mathbb{E}}
\newcommand{\bX}{\mathbf{X}}
\newcommand{\bx}{\mathbf{x}}
\newcommand{\bW}{\mathbf{W}}
\newcommand{\bw}{\mathbf{w}}
\newcommand{\bs}{\mathbf{s}}
\newcommand{\bc}{\mathbf{c}}
\newcommand{\bI}{\mathbf{I}}
\newcommand{\bZero}{\mathbf{0}}
\newcommand{\bM}{\mathbf{M}}
\newcommand{\bo}{\mathbf{o}}
\newcommand{\bO}{\mathbf{O}}
\newcommand{\SoftPlus}{\textrm{SoftPlus}}
\newcommand{\Sigmoid}{\textrm{Sigmoid}}
\newcommand{\InvSoftMax}{\textrm{InvSoftMax}}
\newcommand{\InvSoftPlus}{\textrm{InvSoftPlus}}
\newcommand{\InvSigmoid}{\textrm{InvSigmoid}}
\newcommand{\InvSoftMax}{\textrm{InvSoftMax}}
\newcommand{\bQ}{\mathbf{Q}}
\newcommand{\bV}{\mathbf{V}}
\newcommand{\bU}{\mathbf{U}}
\newcommand{\bu}{\mathbf{u}}
\newcommand{\bp}{\textbf{p}}
\newcommand{\bP}{\textbf{P}}
\newcommand{\ATE}{\textsf{ATE}}
\newcommand{\bepsilon}{\bm{\epsilon}}
\newcommand{\boldeta}{\bm{\eta}}
\newcommand{\Dist}[1]{\textrm{#1}}
\newcommand{\Evolve}[1]{\frac{\textrm{d}#1}{\textrm{d}t}}
\newcommand{\Observe}[1]{\textrm{Observed data }#1}
\newcommand{\bpi}{\boldsymbol{\pi}}
\newcommand{\bPi}{\boldsymbol{\Pi}}
\newcommand{\gst}{g_{\textrm{s.t.}}}
\newcommand{\balpha}{\boldsymbol{\alpha}}
\newcommand{\btheta}{\boldsymbol{\theta}}
\newcommand{\bTheta}{\boldsymbol{\Theta}}
\newcommand{\bgamma}{\boldsymbol{\gamma}}
\newcommand{\ob}{^\textrm{Obs}}
\newcommand{\bv}{\mathbf{v}}
\newcommand{\bC}{\mathbf{C}}
\newcommand{\bz}{\mathbf{z}}
\newcommand{\bZ}{\mathbf{Z}}
\newcommand{\bY}{\mathbf{Y}}
\newcommand{\by}{\mathbf{y}}
\newcommand{\dd}{\textrm{d}}
\newcommand{\bT}{\mathbf{T}}
\newcommand{\bt}{\mathbf{t}}

\newcommand{\bA}{\bm{A}}
\newcommand{\ba}{\bm{a}}
\newcommand{\bh}{\bm{h}}
\newcommand{\bD}{\bm{D}}
\newcommand{\bd}{\bm{d}}
\newcommand{\bd}{\bm{d}}
\newcommand{\cI}{\mathcal{I}}
\newcommand{\cL}{\mathcal{L}}
\newcommand{\cU}{\mathcal{U}}
\newcommand{\cW}{\mathcal{W}}
\newcommand{\cV}{\mathcal{V}}
\newcommand{\cZ}{\mathcal{Z}}
\newcommand{\cD}{\mathcal{D}}
\newcommand{\oD}{\overline{D}}
\newcommand{\oY}{\overline{Y}}
\newcommand{\bone}{\mathbf{1}}

\newcommand{\bbeta}{\boldsymbol{\beta}}
\newcommand{\bsigma}{\boldsymbol{\sigma}}
\newcommand{\blambda}{\boldsymbol{\lambda}}
\newcommand{\bphi}{\boldsymbol{\phi}}
\newcommand{\bpsi}{\boldsymbol{\psi}}

% ==Cross Referencing Different Docs
\usepackage{xr}
%\externaldocument{causalprob-app}

\begin{document} 

\newcommand{\tit}{The Bayesian Modeling of Epidemiological Systems: Combining Machine Learning and Prior Science}
%
%%%%%%%%%%%%%%%%%%%%%%%%%%%%%%%%%%%%%%%%%%%%%%%%%%%%%%%%%%%%%%%%%%%%%%%%%
%
\spacingset{1.25}

\if0\blind

{\title{\bf\tit\thanks{Thank you to X, Y, and Z for helpful comments on previous versions of this paper.}}
  \author{Author A\thanks{Assistant Professor, Department of
      Government, University of Texas at Austin.
158 W 21st, Austin, TX 78712.
      Email:
      \href{mailto:connor.jerzak@austin.utexas.edu}{connor.jerzak@austin.utexas.edu} URL:
      \href{https://www.connorjerzak.com}{ConnorJerzak.com}}
\and
Authors

  \date{
    \today
  }
}
}

\fi


\if1\blind
\title{\bf \tit}
\fi

\maketitle

\pdfbookmark[1]{Title Page}{Title Page}

\thispagestyle{empty}
\setcounter{page}{0}
         
\begin{abstract}
\noindent Abstract here 
\vspace{.25in}
\\ \noindent {\bf Keywords:}  Epidemiological modeling; Machine learning; COVID-19; Bayesian methods  \end{abstract}


%%%%%%%%%%%%%%%%%%%%%%%%%%%%%%%%%%%%%%%%%%%%%%%%%%%

\clearpage
\spacingset{1.83}
\singlespacing 
\tableofcontents

\section{Introduction} 
Let $\gst(a)$ return the $q$ that solves 
\begin{align*}
X &\sim \textrm{Distribution}(q)
\\\textrm{s.t. } \mathbb{E}[X] &= a.
\end{align*}
Why do we need something like this? Because $\E[f(X)] \ne \E[X]$, and the way we impose constraints on the parameters (e.g., to be positive) act like $f(\cdot)$'s.
\\ \vspace{0.1cm}
\\ \Verb|%%%START BAYES ODE%%%|
%%%START BAYES ODE%%%

%%%Priors 
\begin{itemize}
\item $N \leftarrow 10000$  % population scaling factor
\item $\InvSoftPlus(\beta_{l}) \sim \Dist{Normal}(\gst(1),1)$ % transmission rate (# of people infected infects per time unit, or more precisely, fraction of population infected by a person times scaling rate ) (?)
\item $\InvSoftPlus(\sigma) \sim \Dist{Normal}(\gst(1/2),0.1)$ % time rate, exposed to infected (1 over latency time)
\item $\InvSoftPlus(\gamma) \sim \Dist{Normal}(\gst(1/1),0.1)$ % time rate, infected to recovered (1 over recovery time)
\item $\InvSoftPlus(\xi) \sim \Dist{Normal}(\gst(1/20),0.3)$ % time rate, recovered to sus (1 over immune time)
\item $\InvSigmoid(\delta) \sim \Dist{Normal}(\gst(0.01),0.0001)$ % mortality rate (assumed to be the same as the birth rate)
\item $s_{l0} \sim \Dist{Normal}(0, 0.1)$ % initial sus frac (multiply by N to get number) 
\item $e_{l0} \sim \Dist{Normal}(0, 0.1)$ % initial exposed frac (multiply by N to get total number) 
\item $i_{l0}\sim \Dist{Normal}(0, 0.1)$ % initial infected frac (multiply by N to get total number) 
\item $r_{l0} \sim \Dist{Normal}(0, 0.1)$  %initial recovered frac (multiply by N to get total number)
  %%  %%  %%  %%  %%
  % evolving terms 
\item $\Evolve{s_l} = -\SoftPlus( \beta_{l} ) \cdot i_l \cdot s_l \cdot \frac{1}{N} +\SoftPlus( \xi ) \cdot r_l$ 
\item $\Evolve{e_l} = \SoftPlus( \beta_{l} ) \cdot  i_l \cdot  s_l \cdot \frac{1}{N} - \SoftPlus( \sigma ) \cdot  e_l$
\item $\Evolve{i_l} =  \SoftPlus( \sigma ) \cdot e_l - \SoftPlus( \gamma ) \cdot  i_l$ % if set too small, values clipped
\item $\Evolve{r_l} =  \SoftPlus( \gamma ) \cdot  i_l - \SoftPlus( \xi ) \cdot r_l$
\item $\Evolve{\sigma}  = \Neural2{ \sigma , \gamma , \xi , \delta }$
\item $\Evolve{\gamma}  = \Neural2{ \sigma , \gamma , \xi , \delta }$
\item $\Evolve{\xi}  = \Neural2{ \sigma , \gamma , \xi , \delta }$
\item $\Evolve{\delta}  = \Neural2{ \sigma , \gamma , \xi , \delta }$
\item $\Evolve{\beta_l} = \Neural1{ p_l , \beta_l , s_l , e_l , i_l , r_l }$
 \item $\Evolve{p_l} = \Neural1{ p_l , \beta_l , s_l , e_l , i_l , r_l }$
\item $\Observe = \Sigmoid( \delta ) \cdot i_l $ % deaths 
\item $\Observe = p_l$
\end{itemize}
%%%END BAYES ODE%%%

\Verb|%%%END BAYES ODE%%%|

\subsection{Additional Information} 
\begin{itemize}
\item Surrogate posterior mean for parameters (e.g., $\sigma$) and initial conditions (e.g., $i_o$) obtained via $f(\mathbf{D})$, where $f$ is a parametric machine learning model  and $\mathbf{D}$ is input data.
\begin{itemize}
\item Current $f(\cdot)$ is a transformer + LSTM + dense neural network 
\end{itemize}
\item For time-dependent parameters (e.g., $\beta_t$): 
\begin{itemize}
\item Cannot easily apply continuous time $\beta_t$ since unclear how prior would be incorporated across continuous time domain 
\item Instead: discrete resolution $\beta_t$ 
\begin{itemize}
\item Example: in a year, predict (1st, 2nd, 3rd, 4th quarter $\beta$ values)
\item Interpolation used between $t$ slices 
\end{itemize}
\end{itemize}
\end{itemize}

\begin{figure}[H]
 \begin{center}
   \includegraphics[width=0.45\linewidth]{Screenshot 2023-01-04 at
     8.54.33 AM.png}
   \includegraphics[width=0.45\linewidth]{Screenshot 2023-01-04 at 8.55.10 AM.png}
\end{center}
\caption{Predicted $\beta_t$'s. }%\label{fig}
\end{figure}



{\sc Prior Calibration}
\begin{itemize}
\item Priors are important 
\item Reasonable priors over parameters may induce unreasonable prior predictive distribution 
\begin{itemize}
\item Complex aggregation of parameters not easy to conceptualize via individual parameter priors 
\item Adding joint correlations in prior usually not done due to a lack of information
\item One way to explore ``effects'' of prior on induced predictive distribution: 
\begin{itemize}
\item Calculate $\E_{\textrm{Prior}}\left[ \sum_{i=1}^N \log(p(Y_i))\right]$;
\begin{itemize}
\item This is the expected likelihood of the data under the prior specification. 
\item Hopefully, the priors are good and $\E_{\textrm{Prior}}[ p(Y_i)]$ is large 
\end{itemize}
\item Calculate $\nabla_{\textrm{Prior Parameters}} \E_{\textrm{Prior}}\left[ \sum_{i=1}^N \log(p(Y_i))\right]$
\begin{itemize}
\item {\it ``Derivative of the expected likelihood, with respect to the parameters specifying the priors'' }
\item This quantifies how sensitive the expected likelihood is to changes in a given prior parameter. 
\item In the example from the real data, we see that the expected likelihood is most sensitive to choice of $\sigma$ 
\item If priors happened to be at their optimal values $\leadsto$ small gradient values 
\item Big gradient values indicate possibility of poor prior misspecification 
\item Calculated via automatic differentiation with re-parameterized sampling (approach also used in variational inference estimation of approximate posterior)
\end{itemize}
\end{itemize}
\end{itemize}
\end{itemize}

\begin{figure}[H]
 \begin{center}
   \includegraphics[width=0.95\linewidth]{Screenshot 2023-01-04 at 8.44.37 AM.png}
\end{center}
\caption{Some of the model specifications. }%\label{fig}
\end{figure}

\begin{figure}[H]
 \begin{center}
   \includegraphics[width=0.95\linewidth]{Screenshot 2023-01-04 at 8.42.51 AM.png}
\end{center}
\caption{Model results on the real data. Prior period +cumulative infections/hospitalizations/deaths, as well as stringency index used in ML model. $\beta_t$ here has a resolution of 6 (6 distinct $\beta_t$'s predicted, interpolations for in-between periods) }%\label{fig}
\end{figure}

\begin{figure}[H]
 \begin{center}
   \includegraphics[width=0.95\linewidth]{Screenshot 2023-01-04 at 8.42.57 AM.png}
\end{center}
\caption{Model results on the real data. Prior period +cumulative infections/hospitalizations/deaths, as well as stringency index used in ML model. }%\label{fig}
\end{figure}

\begin{figure}[H]
 \begin{center}
   \includegraphics[width=0.95\linewidth]{Screenshot 2023-01-04 at 8.43.05 AM.png}
\end{center}
\caption{Model results on the real data. Prior period +cumulative infections/hospitalizations/deaths, as well as stringency index used in ML model. }%\label{fig}
\end{figure}

\begin{figure}[H]
 \begin{center}
   \includegraphics[width=0.95\linewidth]{Screenshot 2023-01-04 at 8.43.28 AM.png}
\end{center}
\caption{Model results on the real data. Prior period +cumulative infections/hospitalizations/deaths, as well as stringency index used in ML model. }%\label{fig}
\end{figure}

\begin{figure}[H]
 \begin{center}
   \includegraphics[width=0.95\linewidth]{Screenshot 2023-01-04 at 8.43.46 AM.png}
\end{center}
\caption{Model results on the real data. Prior period +cumulative infections/hospitalizations/deaths, as well as stringency index used in ML model. }%\label{fig}
\end{figure}

\begin{figure}[H]
 \begin{center}
   \includegraphics[width=0.95\linewidth]{Screenshot 2023-01-04 at 12.22.08 PM.png}
\end{center}
\caption{Gradient analysis for prior mean terms. }%\label{fig}
\end{figure}

\begin{figure}[H]
 \begin{center}
   \includegraphics[width=0.95\linewidth]{Screenshot 2023-01-04 at 1.02.37 PM.png}
\end{center}
\caption{Gradient analysis for prior variance-covariance terms. }%\label{fig}
\end{figure}

\end{document}

%%% Local Variables:
%%% mode: latex
%%% TeX-master: t
%%% End:
